\section{Fundamental group and integral homology}\label{sec:topology}

This section proves Propositions~\ref{prop:degree-two-mv} and
\ref{prop:integral-topology}.
We adapt the deformation-retraction and Mayer--Vietoris calculations for
$\Lambda$ in
\citep[Sections~7.1--7.6 and Appendix~A]{Alpoge2026}.  For $\Lambda^+$,
the cusp fibre has two components and its boundary cohomology contains
additional classes of order two.  Integral coefficients are therefore
essential.
Throughout this section, all matrices act on column vectors.
On the homology lattice of a smooth fibre we continue to use the basis
$\Lambda^+=\mathbb Z\langle \widehat\gamma,u,w,d\rangle$,
and on its dual the basis
$V^+=\mathbb Z\langle\gamma,u^\vee,w^\vee,\Delta\rangle$, dual to
$(\widehat\gamma,u,w,d)$.  We use the contragredient convention
$T_j(\chi)=\chi\circ A_j^{-1}$, or equivalently
$T_j=(A_j^{-1})^{\mathsf T}$, for $\chi\in V^+$.
This is the action used in the cohomological Wang and Cartan--Leray spectral
sequences below.
For $j=0$, this is the cohomological monodromy associated with the clockwise
cusp meridian used below.  The cusp monodromy
$M_0:=A_0=(A_1A_2)^{-1}$ is
\begin{equation}\label{eq:cusp-monodromy-homology}
 M_0\widehat\gamma=\widehat\gamma-2d,\qquad
 M_0u=u+w,\qquad M_0w=w,\qquad M_0d=d.
\end{equation}
Under this convention, the action on the dual lattice is
\begin{equation}\label{eq:cusp-monodromy-cohomology}
 T_0\gamma=\gamma,\qquad T_0u^\vee=u^\vee,\qquad
 T_0w^\vee=w^\vee-u^\vee,\qquad T_0\Delta=\Delta+2\gamma.
\end{equation}% display-paragraph-boundary: intentional new paragraph follows

Choose pairwise disjoint closed discs about the three exceptional points of
the base.  On the cusp disc write $t=t_c$ for the local coordinate from
Section~\ref{sec:analytic-construction}, so that $p_0$ is given by $t=0$.
For a sufficiently small $\eta>0$, let
$A=f^{-1}(\{|t|\leq\eta\})$ be the closed cusp neighbourhood, let
$Y_{\mathrm{fin}}$ be the inverse
image of the complementary closed disc, and put
$M=A\cap Y_{\mathrm{fin}}=\partial A=\partial Y_{\mathrm{fin}}$.

\subsection{The cusp neighbourhood and its boundary}

At the cusp let
$K=\mathbb Z\langle e_1,e_2\rangle
=\mathbb Z\langle w,d\rangle$ and
$\Gamma=\mathbb Z\langle \widehat\gamma,u\rangle$,
and put $\Gamma'=2\Gamma$ and
$G=\Gamma/\Gamma'\cong(\mathbb Z/2)^2$.  This is the deck group of the
intermediate $\Gamma'$-quotient used below.  The cocharacter-translation
map is the injection
\begin{equation}\label{eq:cusp-B-marking}
 B_0^+:\Gamma\longrightarrow K,
 \qquad B_0^+(\widehat\gamma)=-2d,\qquad B_0^+(u)=w.
\end{equation}
Its image is
$L=\mathbb Zw\oplus\mathbb Z(2d)$.  Write
$T_B=K_{\mathbb R}/L$ and $T_f=K_{\mathbb R}/K$.
The real-linear extension of $B_0^+$ is an isomorphism
$\Gamma_{\mathbb R}\to K_{\mathbb R}$ carrying $\Gamma$ onto $L$, and
hence identifies $\Gamma_{\mathbb R}/\Gamma$ with $T_B$.
The chosen homology marking identifies a nearby real four-torus with
$F=T_B\times T_f$.  The
clockwise cusp meridian is denoted by $a_0$; this convention agrees with
$g_0=(g_1g_2)^{-1}$ and with the logarithmic coordinate
$t=e^{2\pi i s}$, for which $s(g_0z)=s(z)-1$.
For a monodromy-invariant fibre one-cycle $c$, write
$\operatorname{sw}_c(a_0)$ for the suspension torus obtained by
transporting $c$ once around this meridian, oriented by
$a_0\wedge c$.

\begin{lemma}[Comparison of the cusp models]
\label{lem:cusp-model-comparison}
Put $C_{\mathrm{lim}}=C^+(0)$.  After shrinking the cusp disc, the families
$C_\rho(t)=(1-\rho)C^+(t)+\rho C_{\mathrm{lim}}$, for
$0\leq\rho\leq1$, and, in real flat coordinates
$a,b\in\mathbb R^2$, the maps
$\zeta=(sB_0^++C_{\mathrm{lim}})a+b\longmapsto
\zeta-\sigma C_{\mathrm{lim}}a$, for $0\leq\sigma\leq1$,
define homotopies of the quotient torus family over the punctured disc.
They conjugate the corresponding deck actions and extend continuously
through the toric central fibre.  Each map preserves every toric face and
sends cosets of the corresponding torus isotropy subgroup to cosets of that
subgroup.  The extensions are $G$-equivariant.  The extension obtained from
$C_\rho(t)$ is the identity on the central fibre.  Consequently, the
restrictions of the
deformation retractions
supplied by
Lemma~\ref{lem:cusp-retraction} are homotopic on a nearby fibre for the actual
period block and for the model $sB_0^+$.
\end{lemma}

\begin{proof}
For the first family put
$B_{t,\rho}=B_0^+-2\pi\Im C_\rho(t)/\log|t|$.
Uniform invertibility of $B_{t,\rho}$, for $0<|t|<\epsilon$ and
$0\leq\rho\leq1$, follows after shrinking $\epsilon$ from the
invertibility of $B_0^+$.  On a nearby fibre every point has real flat
coordinates
$
 \zeta=(sB_0^++C^+(t))a+b
$
with $(a,b)\in\mathbb R^2\oplus\mathbb R^2$, unique modulo
$\mathbb Z^2\oplus\mathbb Z^2$ in the marked real torus.  Replacing
$(a,b)$ by another integral representative changes the image by the
corresponding period for $C_\rho(t)$.  Hence
\begin{equation}\label{eq:first-cusp-comparison}
 (sB_0^++C^+(t))a+b\longmapsto
 (sB_0^++C_\rho(t))a+b
\end{equation}
is independent of the representative and conjugates the full lattice and
deck actions.

Choose $a$ in a fixed compact fundamental parallelogram.  Since
$C_\rho(t)-C^+(t)=O(t)$, the torus multiplier in
\eqref{eq:first-cusp-comparison} tends uniformly to one as $t\to0$.  Its
logarithmic modulus divided by $\log|t|$ tends to
zero, so it changes neither the limiting fan face nor the corresponding
torus isotropy subgroup.  In every unimodular toric chart it therefore extends continuously
as the identity on each central orbit.  The extensions agree on chart
overlaps because \eqref{eq:first-cusp-comparison} is already a single map on the
dense quotient and the toric space is Hausdorff.

For the constant twist, put
\[
 m_\sigma(y)=\exp\!\left(
 -2\pi i\sigma C_{\mathrm{lim}}
 \bigl((B_0^+\otimes\mathbb R)^{-1}y\bigr)\right)
 \quad(y\in K_{\mathbb R}).
\]
The same deformation argument applies after retaining the finite-index
sublattice $L$; compare \citep[Lemmas~7.5--7.10]{Alpoge2026}.  On the
$K_{\mathbb R}$-cover of $T_B$, where $y=B_0^+a$, one has
\begin{equation}\label{eq:constant-comparison-equivariance}
 m_\sigma(y+B_0^+\lambda)e^{2\pi iC_{\mathrm{lim}}\lambda}
 =e^{2\pi i(1-\sigma)C_{\mathrm{lim}}\lambda}m_\sigma(y),
 \qquad \lambda\in\Gamma.
\end{equation}
Thus multiplication by $m_\sigma(y)$ conjugates the
$C_{\mathrm{lim}}$-twisted deck action to the action with twist
$(1-\sigma)C_{\mathrm{lim}}$, and commutes with the $K$-periods.

In the polar model of a toric chart, the positive part of $m_\sigma(y)$
stays in the same face and its compact part translates the phase.  If
$T_P\subset T_f$ is the isotropy subgroup of a face $P$, this
translation induces a translation of $T_f/T_P$; see
\citep[Proposition~3.1]{NakayamaOgus2010}.  Since $m_\sigma$ is continuous
and bounded on every closed dual cell, these maps extend over the toric
boundary and agree on common faces.  Their logarithmic displacement is
$o(|\log|t||)$, so the extension has the same limiting face as the
map on the punctured family.  Equation
\eqref{eq:constant-comparison-equivariance} also gives $G$-equivariance.
Taking $\rho=\sigma=1$ reduces the family to the
model with period block $sB_0^+$.
\end{proof}

\begin{lemma}[Deformation retraction of the cusp neighbourhood]\label{lem:cusp-retraction}
Let $W=f^{-1}(p_0)\subset A$ be the two-component central fibre determined
by the periodic triangulation $\mathcal T$.  For
a sufficiently small closed cusp disc, the following assertions hold.
\begin{enumerate}
\item The inclusion $W\hookrightarrow A$ is a strong deformation
retract.
\item The restriction of the deformation retraction to a nearby fibre is
homotopic to the collapse map
$c:T_B\times T_f\to W$.
Over the interior of a dual two-cell it leaves $T_f$ unchanged; over the
dual edge corresponding to an edge of $\mathcal T$ with primitive
direction $n$, it quotients by the circle $S_n\subset T_f$ generated by
$n$; over a dual vertex it collapses all of $T_f$.
\item The inclusion of the boundary induces
\begin{equation}\label{eq:cusp-pi1-map}
 \pi_1(M)\longrightarrow\pi_1(A)=\mathbb Z\langle \widehat\gamma,u\rangle,
 \qquad
 \widehat\gamma\mapsto\widehat\gamma,\quad
 u\mapsto u,\quad w,d,a_0\mapsto1.
\end{equation}
\item The suspension tori $\xi_w=\operatorname{sw}_{w}(a_0)$ and
$\xi_d=\operatorname{sw}_{d}(a_0)$
bound embedded solid tori in $A$.
\end{enumerate}
\end{lemma}

\begin{proof}
\smallskip
\noindent\emph{A semistable finite cover and equivariant shrinking.}
The finite cover allows the retraction to be constructed on the closed
neighbourhood while retaining the map induced by its boundary inclusion.
In the fan of cones over $\mathcal T$, pass
from $\Gamma$ to $\Gamma'=2\Gamma$.  Its translation image is
$L'=2L$, and the quotient deck group is $G$.  Every difference of two
vertices
of a simplex of $\mathcal T$ is primitive in $K$, whereas $L'\subset2K$.
Thus no edge or triangle identifies two of its vertices modulo $L'$.
The central fibre of the $\Gamma'$-quotient is consequently a reduced
simple normal crossings divisor, with local equation
$t=z_1\cdots z_r$, where $1\le r\le3$.
The total space is smooth, the map is proper and flat (locally, $t$ is a
non-zero-divisor in a regular local ring), and it is smooth over the
punctured disc.  The $G$-action is free even on the double and
triple strata: if $\bar\lambda\in G$ fixed $[x]$, then
$\Psi_\lambda x=\Psi_\mu x$ for some $\mu\in\Gamma'$; freeness of the
original $\Gamma$-action gives $\lambda=\mu$.

The finite cover is therefore a semistable degeneration of the kind used by
Usui: its total space is smooth, the map to the disc is proper and flat, the
central fibre is reduced with simple normal crossings, and the punctured
fibres are smooth.  We take the inverse image of a closed subdisc preserved
by the radial action.

We apply Usui's monoid action and shrinking map
\citep[Theorem~5.2 and equations (5.2.5.2), (5.2.5.4), and (5.2.12)]{Usui2001}.
We make the auxiliary choices in Usui's construction $G$-equivariant as
follows.  Index the
irreducible components upstairs by a finite set $J$, which $G$ permutes.
Choose a finite $G$-stable system of standard-coordinate charts, each with
trivial stabilizer, and transport coordinates and branches around every
$G$-orbit.  Run the descending inductions of Usui's
Propositions~3.2 and~4.3 simultaneously on the $G$-orbits of subsets of
$J$.  Whenever an indexed partition or cutoff family
$\{p_U\}$ occurs, replace it by
$p_U^G(x)=|G|^{-1}\sum_{h\in G}p_{h^{-1}U}(h^{-1}x)$.
Then $p_{gU}^G(gx)=p_U^G(x)$; subordination, closed-support conditions,
and the fibre-constancy conditions (3.2.22), (3.2.29), and (4.3.4) are
preserved because they are affine and are permuted by $G$.  Induction in
the defining formulas therefore gives $\pi_{gI}\circ g=g\circ\pi_I$ and
$z_{gi}(gy)=z_i(y)$
as systems of branches.  Thus the normal projections and global equations
are $G$-equivariant.

All central multiplicities are one.  Coordinate permutations preserve the
cube, its face projections, and hence Usui's action $R$ on hyperbolic
polar coordinates.  Write $z_i=r_iu_i$, and let $\xi$ be a point of the
logarithmic fibre over $y=\tau_X(\xi)$.  For the radial submonoid,
equation (5.2.12) reads
\begin{equation}\label{eq:usui-radial-phases}
 r_i((s,1)\cdot\xi)=R(s,y)_i,
 \qquad u_i((s,1)\cdot\xi)=u_i(\xi).
\end{equation}
For $(s,1)$, the phase coordinates in
\eqref{eq:usui-radial-phases} are unchanged.  Since $G$ permutes the chosen
coordinates, the radial action is strictly $G$-equivariant.  The homotopy
$H'(x,\rho)=((1-\rho),1)\cdot x$ stays in this closed neighbourhood by (5.2.5.2),
lands in the central fibre at $\rho=1$, and fixes that fibre pointwise by
(5.2.5.4) and \eqref{eq:usui-radial-phases}.  Let $q:A'\to A=A'/G$ be the
finite quotient of this neighbourhood.  Since $q\times\mathrm{id}_{[0,1]}$ is
a quotient map, there is a unique continuous homotopy
$H:A\times[0,1]\to A$ satisfying
$H(q(x),\rho)=q\bigl(H'(x,\rho)\bigr)$.
It fixes $W$ pointwise and has image in $W$ at $\rho=1$, so it is a
strong deformation retraction, proving the first assertion.

\smallskip
\noindent\emph{The collapse map on a nearby fibre.}
The polar-coordinate description gives the following model for the descended
central fibre.  For $b\in T_B$, let $P_b$ be the face of the
triangulation $\mathcal T$ whose open dual cell contains $b$.  If
$P_b$ has vertices $v_0,\ldots,v_r$, let $T_b\subset T_f$ be the
subtorus whose cocharacter lattice is generated by
$v_1-v_0,\ldots,v_r-v_0$.  Define
$(b,x)\sim(b',x')$ precisely when $b=b'$ and $x'-x\in T_b$.  Then
\begin{equation}\label{eq:cusp-collapse-quotient}
 W=(T_B\times T_f)/{\sim}.
\end{equation}
Thus $T_b$ is trivial when $P_b$ is a vertex, is $S_n$ when $P_b$
is an edge of primitive direction $n$, and is all of $T_f$ when
$P_b$ is a unimodular triangle.  These are the three cases in
Lemma~\ref{lem:cusp-retraction}.  The same polar description is the local
model of \citep[Proposition~3.1 and Theorems~3.7, 5.1]{NakayamaOgus2010}.

To identify this map, work first over a sufficiently small positive real
value of $t$.  In the chart belonging to
$P=[v_0,\ldots,v_r]$, write $t=z_0\cdots z_r$ and
$z_i=r_iu_i$.  Equation~\eqref{eq:usui-radial-phases} says that the
radial action changes only the radii.  At its endpoint, projection from the
logarithmic chart to the ordinary toric chart forgets precisely the phases
of the vanishing coordinates.  On the fixed-$t$ fibre their product phase
is fixed, so the endpoint fibre is
\[
 \ker\bigl((S^1)^{r+1}\longrightarrow S^1\bigr)\cong(S^1)^r,
 \qquad
 \mathbb Z\langle(v_i,1):0\leq i\leq r\rangle\cap(K\oplus0)
 =\mathbb Z\langle(v_i-v_0,0):1\leq i\leq r\rangle.
\]
The nonvanishing-coordinate phases are retained.  The description is
invariant under monomial coordinate changes; the local descriptions agree on
overlaps by \citep[p.~30, Claim~2 following (5.2.12)]{Usui2001}.  The radial endpoint
induces a map
$h:T_B\to T_B$ preserving every closed dual cell.  Indeed, a closed dual
cell is characterized in a toric chart by the set of coordinates allowed to
vanish.  Usui's radial formula changes only the corresponding radii and its
face projections commute on inclusions; it neither introduces a coordinate
 outside that face nor changes the formula on a common face.  Hence
$h(\overline P)\subset\overline P$ for every dual cell $P$, including
its lower-dimensional boundary.  In particular, $h$ fixes every dual
vertex.  Let $p:K_{\mathbb R}\to T_B=K_{\mathbb R}/L$ be the universal
covering.  Fix a lift $\widetilde v$ of a dual vertex, and let $\widetilde h$
be the lift of $h\circ p$ fixing it.

On the one-skeleton, $h$
maps each edge into itself relative to its endpoints, so $h_*$ is the
identity on $\pi_1(T_B)$.  It follows that
$\widetilde h(y+\ell)=\widetilde h(y)+\ell$ for
$y\in K_{\mathbb R}$ and $\ell\in L$.
Lifting successively across adjacent cells now shows that
$\widetilde h$ maps each closed lifted cell into itself.
For $0\leq u\leq1$, the explicit interpolation
$\widetilde h_u(y)=(1-u)\widetilde h(y)+uy$
is still $L$-equivariant and stays in the same lifted cell, since that cell
is convex.  This translation-equivariant interpolation is compatible on common faces;
because the isotropy subgroups above are constant on open faces and only
increase on their boundaries, it lifts to a homotopy of the quotient maps.
Thus the restriction of the shrinking map for the model with period block
$sB_0^+$ is homotopic to $c$.

Lemma~\ref{lem:cusp-model-comparison} now removes the holomorphic and constant
deck twists without changing the torus isotropy subgroups.  Hence the
restriction of the actual deformation retraction to $F$ is homotopic to
the collapse map $c$.

\smallskip
\noindent\emph{The fundamental group and the boundary inclusion.}
Projection of \eqref{eq:cusp-collapse-quotient} to $T_B$ has the section
$b\mapsto[b,1]$.  Every loop in $T_f$ can be moved to a dual vertex,
where all of $T_f$ is collapsed.  Cellular van Kampen therefore gives
$\pi_1(A)=\pi_1(T_B)=L$.
Under the identification in \eqref{eq:cusp-B-marking}, its generators
correspond to $\widehat\gamma,u$; the loops $w,d$ in $T_f$ are
null-homotopic.

For the meridian and the suspension tori, use the chart of the
triangle $[0,e_1,e_2]$, use $t=z_0z_1z_2$, $x_1=z_1$, and
$x_2=z_2$.
Here the coordinate meridian is counterclockwise, whereas $a_0$ is
clockwise.  The three coordinate discs are
\[
 \begin{split}
 E_0&=\{z_1=z_2=1,\ |z_0|\le\eta\},\\
 E_w&=\{z_0=\eta,\ z_2=1,\ |z_1|\le1\},\\
 E_d&=\{z_0=\eta,\ z_1=1,\ |z_2|\le1\}.
 \end{split}
\]
Their counterclockwise boundary classes in $H_1(M;\mathbb Z)$ are
\begin{center}
\begin{tabular}{c@{\qquad}c}
\toprule
disc & class in $H_1(M;\mathbb Z)$\\
\midrule
$E_0$ & $-a_0$\\
$E_w$ & $-a_0+w$\\
$E_d$ & $-a_0+d$\\
\bottomrule
\end{tabular}
\end{center}
These discs prove that $a_0,w,d$ die, and hence establish
\eqref{eq:cusp-pi1-map}.

\smallskip
\noindent\emph{The suspension tori.}
Set $V_w=\{z_2=1,\ |z_1|=1,\ |z_0|\le\eta\}$ and
$V_d=\{z_1=1,\ |z_2|=1,\ |z_0|\le\eta\}$.
Each is an embedded solid torus whose boundary lies in $M$.
For $V_w$, if $\phi=\arg z_0$ and $\psi=\arg z_1$, then on the
boundary $(\arg t,\arg x_1,\arg x_2)=(\phi+\psi,\psi,0)$.
The determinant-one change $(\alpha,\beta)=(\phi+\psi,\psi)$ gives
$[\partial V_w]=(-a_0)\wedge(-a_0+w)
=-a_0\wedge w=\pm[\xi_w]$.
The same calculation gives
$[\partial V_d]=-a_0\wedge d=\pm[\xi_d]$.
All displayed coordinate sets project injectively.  Two points of
$E_0,V_w,V_d$ that differ by a deck transformation have the same normalized
logarithmic coordinate $y$, so \eqref{eq:an-log-coordinate-translation} gives
$B_t\lambda=0$.  The matrix $B_t$ tends to $B_0^+$ and is invertible on the
chosen cusp disc, and hence
$\lambda=0$.  On $E_w$ and $E_d$, respectively, the identities
$t=\eta z_1$ and $t=\eta z_2$ first force equality of the two points;
freeness then forces the deck element to be trivial.  On the central
strata, $D_v^\circ$ is sent to $D_{v+B_0^+\lambda}^\circ$, and
injectivity of $B_0^+$ again gives $\lambda=0$.  Since the quotient is
locally biholomorphic, the compact injective images are embedded discs and
solid tori downstairs.
\end{proof}

The cokernel term in the cohomological Wang sequence is already determined by
\eqref{eq:cusp-monodromy-cohomology}:
$C_0:=\operatorname{coker}(T_0-I:V^+\to V^+)
=\mathbb Z[w^\vee]\oplus\mathbb Z[\Delta]
\oplus(\mathbb Z/2)[\gamma]$.
The solid-torus bounds have a cohomological consequence.  Pairing with
$\xi_w$ and $\xi_d$ defines two homomorphisms
$H^2(M;\mathbb Z)\to\mathbb Z$.  On the injected copy of $C_0$ in the
Wang sequence, these homomorphisms extract, with the chosen orientations,
the $[w^\vee]$- and $[\Delta]$-coefficients, respectively; they vanish on
every monodromy-invariant two-cocycle on the fibre.  Since both suspension
tori bound in $A$, the two evaluations vanish on every class restricted from
$H^2(A;\mathbb Z)$.  This formulation does not require a splitting of the
Wang sequence.

\subsection{Homology of the cusp fibre}

The calculation below follows the cellular calculation in
\citep[Appendix~A]{Alpoge2026}, using the quotient dual graph for the
index-two lattice.  The middle panel of
Figure~\ref{fig:cusp-combinatorics} records the dual graph and its translation
labels.
Modulo horizontal translation, the lower and upper triangles at height
$j\in\mathbb Z/2$ give four dual zero-cells
$\mathsf L_j,\mathsf U_j$.  Orient the six dual edges by
$d_j,\upsilon_j:\mathsf L_j\to\mathsf U_j$ and
$h_j:\mathsf L_j\to\mathsf U_{j-1}$.
In the edge order $(d_0,\upsilon_0,h_0,d_1,\upsilon_1,h_1)$, their
translation labels are
$(0,-e_1,-2e_2,0,-e_1,0)$.  There are two dual two-cells
$P_0,P_1$, with $\partial P_0=-d_0+\upsilon_0+d_1-\upsilon_1$ and
$\partial P_1=-\partial P_0$.
Thus the ordinary cellular differentials of the dual torus are
\begin{equation}\label{eq:dual-torus-differentials}
 \partial_1=
 \begin{pmatrix}
 -1&-1&-1& 0& 0& 0\\
  1& 1& 0& 0& 0& 1\\
  0& 0& 0&-1&-1&-1\\
  0& 0& 1& 1& 1& 0
 \end{pmatrix},
 \qquad
 \partial_2=
 \begin{pmatrix}
 -1& 1\\ 1&-1\\ 0&0\\ 1&-1\\-1&1\\0&0
 \end{pmatrix}.
\end{equation}
The nonzero Smith invariants of $\partial_1$ are $1,1,1$, and that of
$\partial_2$ is $1$.  Moreover $\operatorname{im}\partial_2$ is the
primitive rank-one subgroup generated by
$-d_0+\upsilon_0+d_1-\upsilon_1$ inside $\ker\partial_1$.  Thus the
bottom row has homology $(\mathbb Z,\mathbb Z^2,\mathbb Z)$; these Smith
invariants also prevent torsion from entering this row of the filtration
spectral sequence.

For a dual cell $\sigma$, write $T_\sigma=T_b$ for $b$ in its interior,
and filter $W$ by the inverse images, under the projection $W\to T_B$, of
the skeleta of the dual complex.  Over $\sigma$, this projection has fibre
$T_f/T_\sigma$, and the filtration spectral sequence has
$E^1_{p,q}=\bigoplus_{\sigma^p}H_q(T_f/T_\sigma;\mathbb Z)$.
\begin{samepage}
Its nonzero terms are
\begin{center}
\begin{tabular}{c@{\quad}ccc}
\toprule
 & $p=0$ & $p=1$ & $p=2$\\
\midrule
$q=2$ & 0 & 0 & $\mathbb Z^2$\\
$q=1$ & 0 & $\mathbb Z^6$ & $K\oplus K$\\
$q=0$ & $\mathbb Z^4$ & $\mathbb Z^6$ & $\mathbb Z^2$\\
\bottomrule
\end{tabular}
\end{center}
\end{samepage}
The bottom differential is \eqref{eq:dual-torus-differentials}.  Put
$n_d=e_2-e_1$, $n_v=e_2$, and $\nu_n(k)=\det(n,k)$, and write
$\nu_d=\nu_{n_d}$ and $\nu_v=\nu_{n_v}$.
The only differential in the $q=1$ row is
\begin{equation*}
 \begin{split}
 \delta:K\oplus K&\longrightarrow\mathbb Z^6,\\
 (k,k')&\longmapsto
 \bigl(-\nu_d(k-k'),\nu_v(k-k'),0,
       \nu_d(k-k'),-\nu_v(k-k'),0\bigr).
 \end{split}
\end{equation*}
The zero horizontal entries result from the two opposite incidences of the
same horizontal edge orbit.  The map
$K\to\mathbb Z^2$, $k\mapsto(\nu_d(k),\nu_v(k))$, is unimodular.
Therefore $\ker\delta=\{(k,k):k\in K\}$, and its cokernel is the free
abelian group
$\mathbb Z\langle\bar h_0,\bar h_1,\bar d,\bar\upsilon\rangle$, where
$\bar d=[d_0]=[d_1]$ and
$\bar\upsilon=[\upsilon_0]=[\upsilon_1]$.
Equivalently, the two nonzero Smith invariants of $\delta$ are both one.
Every possible higher differential has zero target.  Consequently there
are no torsion or additive extension ambiguities, and
\begin{equation}\label{eq:cusp-homology}
 H_k(A;\mathbb Z)=H_k(W;\mathbb Z)=
 (\mathbb Z,\mathbb Z^2,\mathbb Z^5,
   \mathbb Z^2,\mathbb Z^2),\qquad 0\le k\le4.
\end{equation}% display-paragraph-boundary: intentional new paragraph follows

In degree two, the two
base cycles are $\beta_1=d_0-\upsilon_0$ and
$\beta_2=d_0+d_1-h_0-h_1$; their translation labels are
$e_1$ and $2e_2$, respectively.  If
$\beta=\sum m_e e$ and $v\in K$, the collapse over an edge gives
$c_*(\beta\otimes v)=\sum_e m_e\nu_{n_e}(v)[C_e]$.
Here $C_e$ is the oriented surviving circle in $T_f/S_{n_e}$ over the
dual edge $e$.
With compatible orientations, the four mixed generators have images
\begin{equation}\label{eq:specialization-table}
\begin{array}{c|cc}
 &e_1&e_2\\ \hline
\beta_1&-\bar d+\bar\upsilon&-\bar d\\
\beta_2&-2\bar d&-2\bar d-\bar h_0-\bar h_1.
\end{array}
\end{equation}
The base orientation class maps to a primitive class $\omega_B$.  The phase
orientation class maps to zero because it may be represented over a dual
vertex, where all of $T_f$ is collapsed.  It follows directly from
\eqref{eq:specialization-table} that
\begin{equation}\label{eq:specialization-image}
 c_*H_2(F;\mathbb Z)
 =\mathbb Z\langle\omega_B,\bar d,\bar\upsilon,
                  \bar h_0+\bar h_1\rangle.
\end{equation}
This subgroup is saturated.  Its cokernel is infinite cyclic, primitively
generated by $\bar h_0$, with $\bar h_1=-\bar h_0$; in particular,
$\bar h_0-\bar h_1$ is twice, rather than once, a primitive generator.
With target basis
$(\omega_B,\bar h_0,\bar h_1,\bar d,\bar\upsilon)$ and source basis
$([T_B],[T_f],\beta_1e_1,\beta_1e_2,\beta_2e_1,\beta_2e_2)$, the
specialization matrix is
\begin{equation*}
 \begin{pmatrix}
 1&0& 0& 0& 0& 0\\
 0&0& 0& 0& 0&-1\\
 0&0& 0& 0& 0&-1\\
 0&0&-1&-1&-2&-2\\
 0&0& 1& 0& 0& 0
 \end{pmatrix}.
\end{equation*}
Its four nonzero Smith invariants are all one, so the image in
\eqref{eq:specialization-image} is saturated.

On cohomology put
\begin{equation*}
 I_0=(\textstyle\bigwedge^2V^+)^{T_0}
 =\mathbb Z\langle
 \gamma u^\vee,\gamma\Delta,u^\vee w^\vee,
 u^\vee\Delta-2\gamma w^\vee\rangle.
\end{equation*}
Via $A\simeq W$, restriction to a nearby fibre is the pullback $c^*$.
Since the relevant homology groups are free, its matrix is the transpose of
the specialization matrix above.  Its image is therefore saturated of rank
four, and its kernel is primitive of rank one.  The image lies in $I_0$ by
monodromy invariance; since both are saturated rank-four sublattices of
$\bigwedge^2V^+$ with the same rational span, they are equal.  Thus
\begin{equation}\label{eq:cusp-H2-splitting}
 0\longrightarrow\mathbb Z
 \longrightarrow H^2(A;\mathbb Z)
 \xrightarrow{\ \operatorname{res}_F\ } I_0\longrightarrow0.
\end{equation}
The sequence splits as groups, but no splitting will be used yet.  Meanwhile
the Wang sequence and \eqref{eq:cusp-monodromy-cohomology} give
\begin{equation}\label{eq:cusp-Wang}
 0\longrightarrow C_0\longrightarrow H^2(M;\mathbb Z)
 \longrightarrow I_0\longrightarrow0,
 \qquad
 C_0=\mathbb Z[w^\vee]\oplus\mathbb Z[\Delta]
       \oplus(\mathbb Z/2)[\gamma].
\end{equation}
Write $\vartheta=[\gamma]$ for the unique nonzero torsion element of $C_0$.

\begin{lemma}[The cusp double curve]
\label{lem:cusp-double-curve}
Let $D$ be either component of the cusp fibre, let
$\nu:S\to D$ be its normalization, and let
$C_{\mathrm{dbl}}\cong\mathbb P^1$ be its double curve.  Number the boundary
curves of $S\cong\dP$ as in Figure~\ref{fig:cusp-combinatorics}, so that
$C_1,C_4$ are the two conductor branches, and put
$C=C_2+C_3+C_5+C_6$.  Then
\begin{equation}\label{eq:normalization-normal-bundle}
 \nu^*(\mathcal O_Y(D)|_D)=\mathcal O_S(-C).
\end{equation}
Moreover, $D^3=0$, $D\cdot C_{\mathrm{dbl}}=-2$, and
$\deg N_{C_{\mathrm{dbl}}/Y}=-2$.
\end{lemma}

\begin{proof}
Let $D'$ be the other cusp component.  Since $Y$ is smooth, both
components are Cartier.  On a holomorphic neighbourhood of the cusp fibre,
$D+D'=\operatorname{div}(t)$.  The pullback of $D'|_D$ to the
normalization is the reduced divisor $C$, which proves
\eqref{eq:normalization-normal-bundle}.  Each $C_i$ is a $(-1)$-curve,
and the only adjacent pairs among the four components of $C$ are
$(C_2,C_3)$ and $(C_5,C_6)$.  Hence
$C^2=-4+2(1+1)=0$.  The projection formula for the finite normalization
therefore gives
$D^3=\int_S c_1(\mathcal O_S(-C))^2=C^2=0$.
Either conductor branch, say $C_1$, meets $C$ in two reduced points and
maps isomorphically to $C_{\mathrm{dbl}}$.  Restricting
\eqref{eq:normalization-normal-bundle} to $C_1$ gives
$D\cdot C_{\mathrm{dbl}}=-2$.  Locally along the double curve the two
branches of $D$ have product equation $z_1z_2=0$.  The product of their
normal lines is both $\det N_{C_{\mathrm{dbl}}/Y}$ and
$\mathcal O_Y(D)|_{C_{\mathrm{dbl}}}$, proving the last equality.
\end{proof}

\begin{proposition}[The class of a cusp component]
\label{prop:cusp-component-class}
There is a topological line bundle $\mathcal E$ on the closed cusp neighbourhood
whose first Chern class $\kappa=c_1(\mathcal E)$ generates the kernel in
\eqref{eq:cusp-H2-splitting} and satisfies
\begin{equation}\label{eq:cusp-component-class}
 [D_0]=2\kappa,\qquad [D_1]=-2\kappa.
\end{equation}
Its restriction to a smooth fibre has zero Chern class, while on the boundary
mapping torus
\begin{equation}\label{eq:kappa-boundary}
 \kappa|_M=\vartheta=[\gamma].
\end{equation}
Thus $\mathcal E$ is topologically trivial on every smooth fibre, although
its restriction to $M$ is not.
\end{proposition}

\begin{proof}
Let $D_{(a,j)}$ denote the height-one ray divisor on the infinite toric
cover and set
\[
 E=\sum_{(a,j)\in\mathbb Z^2}
   -\left\lceil\frac j2\right\rceil D_{(a,j)}.
\]
The fan is locally finite, and only the finitely many ray divisors belonging
to a fixed cone meet its affine chart.  Hence this sum is a locally finite
torus-invariant Weil divisor.  The toric cover is smooth, so its local rings
are factorial and $E$ is Cartier on every chart.  The ratios of the local
equations on overlaps are units and therefore define $\mathcal O(E)$.

Let $D_{\mathrm{ev}}$ be the sum over even $j$, and denote by $D_0$
and $D_1$ the components obtained from the divisors with even and odd
values of $j$, respectively.  For the character
$m=(0,1,1)$, evaluation on every height-one ray gives
\begin{equation}\label{eq:half-divisor-identity}
 2E=D_{\mathrm{ev}}-\operatorname{div}(\chi^m).
\end{equation}
Now write $\lambda=r_\lambda\widehat\gamma+s_\lambda u$.  Its height-one translation
is $(s_\lambda,-2r_\lambda)$, and the torus-translation factor fixes every
invariant ray divisor.  Coefficientwise comparison therefore gives
$\Psi_\lambda^*E=E+r_\lambda\operatorname{div}(t)$.
Multiplication by $t^{r_\lambda}$ gives the regular isomorphism
$\Psi_\lambda^*\mathcal O(E)
=\mathcal O(E+r_\lambda\operatorname{div}(t))
\longrightarrow\mathcal O(E)$.
This remains regular when $r_\lambda<0$: the pole of the rational
multiplier is exactly cancelled by the change of the divisor sheaf in the
source.  Since $r_{\lambda+\mu}=r_\lambda+r_\mu$ and
$\Psi_\lambda^*t=t$, these~isomorphisms satisfy the strict cocycle
condition.  They therefore descend $\mathcal O(E)$ to a line bundle
$\mathcal E$ on the cusp neighbourhood.

Pulling the character $\chi^m$ through the holomorphic torus translation
introduces a nowhere-vanishing holomorphic unit $u_\lambda(t)$.  The deck
law makes $\{u_\lambda\}$ a multiplicative cocycle.  On a sufficiently
small disc, choose holomorphic logarithms on the two generators of
$\Gamma$ and extend them additively.  The cocycles
$\exp(\sigma\log u_\lambda)$, $0\leq\sigma\leq1$, give a homotopy from
this unit cocycle to the trivial one.  Thus it contributes no integral first
Chern class, and the descended form of
\eqref{eq:half-divisor-identity} yields $[D_0]=2\kappa$.  Since
$D_0+D_1=\operatorname{div}(t)$ on the cusp neighbourhood, the second equality in
\eqref{eq:cusp-component-class} follows.  Restriction to a nearby fibre
leaves only constant character multipliers, so $c_1(\mathcal E|_F)=0$.

On $t=\eta e^{2\pi i\theta}$, the positive real factor
$\eta^{r_\lambda}$ is homotopically trivial and the descent multiplier is
$e^{2\pi i\theta r_\lambda}=e^{2\pi i\theta\gamma(\lambda)}$.
In the Wang sequence \eqref{eq:cusp-Wang}, this cocycle represents
$[\gamma]\in C_0$, and hence proves
\eqref{eq:kappa-boundary}.

Finally, let $C_{\mathrm{dbl}}$ be the double curve obtained by identifying
the two horizontal boundary curves in the normalization of $D_0$.
Lemma~\ref{lem:cusp-double-curve} gives
$D_0\cdot C_{\mathrm{dbl}}=-2$.  Therefore
\eqref{eq:cusp-component-class} gives
$\kappa(C_{\mathrm{dbl}})=-1$.  Thus $\kappa$ is primitive.  Since its
restriction to $F$ vanishes, it generates the kernel in
\eqref{eq:cusp-H2-splitting}.
\end{proof}

\subsection{The two multiple fibres}

The image of $Y_{\mathrm{fin}}$ in the base is a closed disc containing
the orbifold points $p_1$ and $p_2$.  Split it into two closed discs
$U_1,U_2$ meeting in a band, with $p_j$ in
$U_j$, and put $Y_j=f^{-1}(U_j)$.  On the orbifold cover of
$U_j$, the homotopy
$R_\rho(s_j,x)=((1-\rho)s_j,x)$, for $0\leq\rho\leq1$,
commutes with
$(s_j,x)\mapsto(\zeta_js_j,A_jx+v_j/m_j)$, where
$\zeta_j=e^{-2\pi i/m_j}$.  It therefore descends to a
strong deformation retraction of $Y_j$ onto its reduced central fibre
$S_j$.  Hence
$Y_j\simeq S_j$.  After trivializing the torus bundle over the overlap
band, $Y_1\cap Y_2\simeq F\times[0,1]\searrow F$.
This retraction is compatible with the two retractions $Y_j\searrow S_j$,
and the induced attaching maps are the coverings
$\pi_j:F\to S_j$.  Thus $Y_{\mathrm{fin}}$ has the homotopy-pushout
model $S_1\cup_F S_2$.  Define
$\eta_1=4u^\vee+2w^\vee+3\Delta$,
$\eta_2=u^\vee+w^\vee+\Delta$, and
$Q=u^\vee w^\vee+3\gamma\Delta$.  Here $T_j$ denotes the
contragredient action fixed at the beginning of this section.
The affine Bieberbach group of the reduced fibre is
\[
 \mathcal G_j=
 \left\langle\Lambda^+,r_j\ \middle|\
 r_j\lambda r_j^{-1}=A_j\lambda,\quad r_j^{m_j}=v_j
 \right\rangle,
 \qquad (m_1,m_2)=(3,4),
\]
and fits into
$1\to\Lambda^+\to\mathcal G_j\to C_{m_j}\to1$.
The spectral-sequence calculation below follows the method of
\citep[Appendix~A, Lemma~A.4]{Alpoge2026}, with $\Lambda^+$ in place of
$\Lambda$; in particular, the extension associated with $S_2$ changes.
For the free affine action on $F$, the Cartan--Leray spectral sequence,
equivalently the Lyndon--Hochschild--Serre sequence of this extension, is
\begin{equation*}
 E_2^{p,q}=H^p\!\left(C_{m_j};\textstyle\bigwedge^qV^+\right)
 \Longrightarrow H^{p+q}(S_j;\mathbb Z).
\end{equation*}
The edge map is the covering pullback $\pi_j^*$.  In degree one, an
invariant character $\chi$ descends exactly when
$\chi(v_j)\equiv0\pmod{m_j}$.  Hence
$\pi_1^*H^1(S_1;\mathbb Z)=
\mathbb Z\langle3\gamma,\eta_1\rangle$ and
$\pi_2^*H^1(S_2;\mathbb Z)=
\mathbb Z\langle4\gamma,\eta_2\rangle$.

We determine the images of $\pi_j^*$ on $H^2$ without using a
transgression formula.  Transfer identifies $H^2(S_j;\mathbb Q)$ with
$(\bigwedge^2V^+_{\mathbb Q})^{T_j}$.  Direct calculation from the matrices
$A_j$ gives the integral invariant lattices
$I_1=\mathbb Z\langle \gamma\eta_1,Q\rangle$ and
$I_2=\mathbb Z\langle \gamma\eta_2,Q\rangle$.
With the orientation
$\gamma u^\vee w^\vee\Delta$, their Gram matrices in the displayed bases
are
\[
 \begin{pmatrix}0&3\\3&6\end{pmatrix},
 \qquad
 \begin{pmatrix}0&1\\1&6\end{pmatrix},
\]
of determinants $-9$ and $-1$, respectively.  Put
$L_j=\pi_j^*H^2(S_j;\mathbb Z)_{\mathrm{free}}\subset I_j$.
The projection formula gives
$\langle\pi_j^*x\smile\pi_j^*y,[F]\rangle
=m_j\langle x\smile y,[S_j]\rangle$.
The intersection form on $H^2(S_j;\mathbb Z)_{\mathrm{free}}$ is
unimodular, so $|\det L_j|=m_j^2$.  It follows that
$[I_1:L_1]=1$ and $[I_2:L_2]=4$, and in particular
$L_1=\mathbb Z\langle \gamma\eta_1,Q\rangle$.

To identify the index-four sublattice $L_2$, put
$k_1=u-d$ and $k_2=w-d$.  At $p_2$, the monodromy and translation
vector satisfy $A_2k_1=k_2$, $A_2k_2=-k_1$, $A_2d=d$, and
$v_2=-\widehat\gamma-3k_1+3k_2$.
Eliminating $\widehat\gamma$ with $r_2^4=v_2$ gives
\begin{equation}\label{eq:order-four-product-model}
 \pi_1(S_2)=
 \bigl(\mathbb Z^2\langle k_1,k_2\rangle
       \rtimes_R\mathbb Z\langle r_2\rangle\bigr)
 \times\mathbb Z\langle d\rangle,
 \qquad R(k_1)=k_2,\quad R(k_2)=-k_1.
\end{equation}
Both $S_2$ and $N_R\times S^1_d$ are aspherical: the former is a free
finite affine quotient of $F$, while the latter has universal cover
$\mathbb R^4$.  The displayed fundamental-group identification therefore
yields $S_2\simeq N_R\times S^1_d$, where $N_R$ is the mapping torus of
the quarter-turn $R$ on the two-torus.  We write an element of the displayed
semidirect product in the normal form $(v,n;m)$, meaning
$v r_2^n d^m$, with $v\in\mathbb Z^2\langle k_1,k_2\rangle$.
In these coordinates
$\widehat\gamma=(-3k_1+3k_2,-4;0)$.
The Wang and K\"unneth sequences give two generators for the free part of
$H^2(N_R\times S^1_d;\mathbb Z)$: the fibre-orientation class on $N_R$
and the product of the mapping-torus base class with the $d$-circle class.
The first pulls back to
$(u^\vee-3\gamma)(w^\vee+3\gamma)=Q-3\gamma\eta_2$.  For the second, the
normal-form expression for $\widehat\gamma$ shows that the base class pulls
back to $-4\gamma$, while the $d$-circle class pulls back to $\eta_2$.
Thus, up to signs, the two pullback classes are
$Q-3\gamma\eta_2$ and $4\gamma\eta_2$.
Consequently
$L_2=\mathbb Z\langle4\gamma\eta_2,Q-3\gamma\eta_2\rangle
=\mathbb Z\langle4\gamma\eta_2,\gamma\eta_2+Q\rangle$, and the pullback lattices are
\begin{equation*}
\begin{array}{c|c|c}
 &\pi_j^*H^1(S_j;\mathbb Z)
 &\pi_j^*H^2(S_j;\mathbb Z)_{\mathrm{free}}\\ \hline
j=1&\mathbb Z\langle3\gamma,\eta_1\rangle
   &L_1=\mathbb Z\langle \gamma\eta_1,Q\rangle\\
j=2&\mathbb Z\langle4\gamma,\eta_2\rangle
   &L_2=\mathbb Z\langle4\gamma\eta_2,\gamma\eta_2+Q\rangle.
\end{array}
\end{equation*}
Comparing coefficients gives
\begin{equation}\label{eq:finite-fibre-lattice-arithmetic}
 L_1\cap L_2=4\mathbb ZQ,
 \qquad
 L_1+L_2=\mathbb Z\langle\gamma\eta_1,\gamma\eta_2,Q\rangle.
\end{equation}
The latter sum is saturated: on the coordinates
$(\gamma w^\vee,\gamma\Delta,u^\vee w^\vee)$, the three displayed generators have a
minor of determinant $-1$.

We first construct the torsion class on $S_2$.  The homological
coinvariants are
$\Lambda^+/(A_2-I)\Lambda^+=\mathbb Z[\widehat\gamma]\oplus\mathbb Z[u]
\oplus(\mathbb Z/2)[d-u]$.  Since the affine relation in these coinvariants
is $4r_2=v_2=-\widehat\gamma$, eliminating $\widehat\gamma$ gives
$H_1(S_2;\mathbb Z)=\mathbb Z[r_2]\oplus\mathbb Z[u]
\oplus(\mathbb Z/2)[d-u]$.  Thus the universal coefficient theorem gives a
unique nonzero torsion class
$\alpha_2\in H^2(S_2;\mathbb Z)$, with
$\operatorname{Tor}H^2(S_2;\mathbb Z)\cong\mathbb Z/2$.

To determine its boundary image, consider the normal line of $S_2$.  The
affine generator $r_2$ acts on this line by the unitary
character $r_2\mapsto e^{-2\pi i/4}$, which is trivial on $\Lambda^+$.
This character has the real lift
$\widetilde\chi(\lambda)=\gamma(\lambda)$ and
$\widetilde\chi(r_2)=-1/4$,
because $\gamma\circ A_2=\gamma$ and
$4\widetilde\chi(r_2)=-1=\gamma(v_2)$.  Thus the normal line is
topologically trivial.  In the exponential sequence
$0\to\mathbb Z\to\mathbb R\to U(1)\to0$, the real character
$\widetilde\chi$ lifts the unitary normal character, so its Bockstein,
which is the first Chern class of the normal line, vanishes.  Let
$p:M_2=\partial Y_2\to S_2$ be its unit-circle bundle.  Its Euler class is
therefore zero, and the Gysin sequence gives
$\ker(p^*:H^q(S_2;\mathbb Z)\to H^q(M_2;\mathbb Z))
=\operatorname{im}({-}\smile c_1(p))=0$.
Thus $p^*$ is injective in every degree.

The relations in the cokernel term of the cohomological Wang sequence are
$6\gamma-u^\vee+w^\vee=0$, $-u^\vee-w^\vee=0$, and
$-6\gamma+2u^\vee=0$.  Consequently,
$\operatorname{coker}(T_2-I|V^+)
=\mathbb Z[\gamma]\oplus\mathbb Z[\Delta]
\oplus(\mathbb Z/2)[u^\vee-3\gamma]$.
In this local cokernel $[\gamma]$ has infinite order, whereas the class
denoted $[\gamma]$ in $C_0$ has order two.  Since $p^*$ is injective,
$p^*\alpha_2$ is represented here by $[u^\vee-3\gamma]$.  This determines
the restriction to $M_2$; the restriction to the common boundary $M$ is
determined below, after gluing, by naturality of the universal-coefficient
$\operatorname{Ext}$-term.

Combining the presentations of the two affine fundamental groups in the
homotopy pushout, the $A_1$-coinvariant relations reduce to
$u=2w$ and $2d=3w$, whereas the $A_2$-coinvariant relations reduce to
$u=w$ and $2(d-u)=0$.  Together they give $u=w=0$ and $2d=0$.
The affine relations then become
$3r_1=\widehat\gamma$ and $4r_2=-\widehat\gamma$.  Hence
\begin{equation}\label{eq:H1-B-presentation}
 \begin{split}
 H_1(Y_{\mathrm{fin}};\mathbb Z)
 &=\frac{\mathbb Z\langle\widehat\gamma,d,r_1,r_2\rangle}
 {\mathbb Z\langle2d,\,3r_1-\widehat\gamma,\,
 4r_2+\widehat\gamma\rangle}\\
 &\cong\mathbb Z\langle\mathfrak s\rangle
   \oplus(\mathbb Z/2)\langle d\rangle,
 \end{split}
\end{equation}
where $\widehat\gamma=12\mathfrak s$, $r_1=4\mathfrak s$, and
$r_2=-3\mathfrak s$.  The cohomology
Mayer--Vietoris segment is
\[
 \begin{split}
 H^1(S_1)\oplus H^1(S_2)&\longrightarrow H^1(F)
 \longrightarrow H^2(Y_{\mathrm{fin}})\\
 &\longrightarrow H^2(S_1)\oplus H^2(S_2)
 \longrightarrow H^2(F).
 \end{split}
\]
The cokernel of its first arrow is free cyclic.  Indeed,
$3\gamma$ and $4\gamma$ generate $\gamma$, while
imposing the relations $\eta_1-3\eta_2=u^\vee-w^\vee$ and $\eta_2=0$ gives
$u^\vee=w^\vee$, $\Delta=-2w^\vee$.  Thus this cokernel is generated primitively by
$[w^\vee]$.  By \eqref{eq:finite-fibre-lattice-arithmetic}, the kernel of
the last arrow on the free summands is $\mathbb Z(4Q)$.  In addition, the
class $\alpha_2$ of order two from $S_2$ lies in the full kernel and
restricts trivially
to the torsion-free group $H^2(F)$.  The universal coefficient theorem and
\eqref{eq:H1-B-presentation} give an injection
$\operatorname{Ext}(H_1(Y_{\mathrm{fin}}),\mathbb Z)\cong\mathbb Z/2$
into $H^2(Y_{\mathrm{fin}};\mathbb Z)$; naturality for
$S_2\hookrightarrow Y_{\mathrm{fin}}$ is expressed by
\[
 \begin{array}{ccc}
 \operatorname{Ext}(H_1(Y_{\mathrm{fin}}),\mathbb Z)
   &\lhook\joinrel\longrightarrow&H^2(Y_{\mathrm{fin}};\mathbb Z)\\
 \Big\downarrow&&\Big\downarrow\\
 \operatorname{Ext}(H_1(S_2),\mathbb Z)
   &\lhook\joinrel\longrightarrow&H^2(S_2;\mathbb Z).
 \end{array}
\]
The left arrow is an isomorphism because
$[d-u]\mapsto[d]$ is an isomorphism on torsion.  Hence the nonzero torsion
class restricts to $\alpha_2$.  In particular,
$H^2(Y_{\mathrm{fin}};\mathbb Z)$ contains an element of order two mapping to
$\alpha_2$.  Hence the extension of the $\mathbb Z/2$-summand by the
connecting $\mathbb Z$-summand splits, and
\begin{equation}\label{eq:H2-B}
 H^2(Y_{\mathrm{fin}};\mathbb Z)=
 \mathbb Z b\oplus\mathbb Z Q_{\mathrm{fin}}
 \oplus(\mathbb Z/2)\alpha.
\end{equation}
Here $b$ is the connecting class represented by $w^\vee$ in
$V^+/(\pi_1^*H^1(S_1)+\pi_2^*H^1(S_2))\cong\mathbb Z[w^\vee]$, the class
$Q_{\mathrm{fin}}$ restricts to $4Q$ on $F$, and $\alpha$ is the
unique nonzero torsion class.  Its existence and uniqueness also follow directly
from \eqref{eq:H1-B-presentation} and universal coefficients.

The summand of order two in \eqref{eq:H2-B} is the
$\operatorname{Ext}$-term in the universal coefficient theorem arising
from $\operatorname{Tor}H_1(Y_{\mathrm{fin}};\mathbb Z)$.
The next lemma concerns the homology group $H_2(Y_{\mathrm{fin}};\mathbb Z)$.

\begin{lemma}[Pushforward under the fourfold cover]
\label{lem:fourfold-pushforward}
For the fourfold covering $\pi_2:F\to S_2$, one has
$H_2(S_2;\mathbb Z)=\mathbb Z^2\oplus
(\mathbb Z/2)\langle z\rangle$ and $\pi_{2*}(u\wedge d)=z$.
Consequently $H_2(Y_{\mathrm{fin}};\mathbb Z)$, and hence
$H^3(Y_{\mathrm{fin}};\mathbb Z)$, is torsion free.
\end{lemma}

\begin{proof}
Retain the notation and product model of
\eqref{eq:order-four-product-model}.
Write $\bar k_1$ for the class of $k_1$ in the coinvariants
$\mathbb Z^2/(R-I)\mathbb Z^2$.  The Wang and K\"unneth sequences
therefore give
$H_2(S_2;\mathbb Z)=\mathbb Z^2\oplus
(\mathbb Z/2)\langle z\rangle$, where
$z=[\bar k_1]\times[d]$.
For comparison, the abelianization of $\pi_1(S_1)$ has
$u=2w$, $2d=3w$, and $\widehat\gamma=3r_1$.  Coprimality of $2$ and $3$ gives
$H_1(S_1;\mathbb Z)\cong\mathbb Z^2$, and Poincar\'e duality gives
torsion-free $H_2(S_1;\mathbb Z)$.
For the fourfold covering, the integral pushforward is
\begin{equation}\label{eq:order-two-pushforward}
 \pi_{2*}(u\wedge d)
 =\pi_{2*}(k_1\wedge d)=z.
\end{equation}
Here $\pi_{2*}$ is induced by the quotient covering, not the transfer, so
no factor of four occurs.  Indeed, under the product model above it sends the
fibre loop $k_1$ to its coinvariant $\bar k_1$ and leaves the $d$-circle
unchanged.  Choose the product cell structure in which $e_{k_1}$ and
$e_d$ are the corresponding oriented one-cells.  The $C_4$-orbit of
$e_{k_1}\times e_d$ consists of the four distinct cells with first
directions $k_1,k_2,-k_1,-k_2$, and hence has trivial stabilizer.  The
quotient characteristic map is
\[
 \begin{array}{ccc}
 D^1\times D^1&
 \xrightarrow{\ \chi_{k_1}\times\chi_d\ }&F\\
 \big\Vert&&\big\downarrow{\pi_2}\\
 D^1\times D^1&
 \xrightarrow{\ \chi_{\bar k_1}\times\chi_d\ }&S_2 .
 \end{array}
\]
On cell interiors $\pi_2$ has degree $+1$ with these orientations, and
the boundary gives the coinvariant relation.  Thus the cellular chain map
sends $k_1\times d$ to $[\bar k_1]\times[d]$, proving
\eqref{eq:order-two-pushforward}.
On the other hand, $u\wedge d$ pairs to zero with each of
$\gamma\eta_1,\gamma\eta_2,Q$, so it lies in the kernel of the free part of
$\Phi_2=(\pi_{1*},-\pi_{2*}):H_2(F)\to H_2(S_1)\oplus H_2(S_2)$
and maps onto all target torsion.  The transpose of the free map has image
$L_1+L_2$, which is saturated by
\eqref{eq:finite-fibre-lattice-arithmetic}; hence the free cokernel is
free.  Let $T=\operatorname{Tor}H_2(S_2)$, and let
$\Phi_{2,\mathrm{free}}$ be the composite of $\Phi_2$ with the
maximal free quotient of its target.  The restriction of $\Phi_2$ to
$\ker\Phi_{2,\mathrm{free}}$ lands in $T$; denote it by $t$.  The
exact sequence
\[
 0\longrightarrow T/t(\ker\Phi_{2,\mathrm{free}})
 \longrightarrow\operatorname{coker}\Phi_2
 \longrightarrow\operatorname{coker}\Phi_{2,\mathrm{free}}
 \longrightarrow0
\]
shows, using \eqref{eq:order-two-pushforward}, that
$\operatorname{coker}\Phi_2$ is free.  Homological
Mayer--Vietoris for $Y_1\cup_F Y_2$ then makes
$H_2(Y_{\mathrm{fin}};\mathbb Z)$ free.  The universal coefficient theorem gives
\begin{equation}\label{eq:H3-B-torsion-free}
 H^3(Y_{\mathrm{fin}};\mathbb Z)\ \text{is torsion free}.
\end{equation}
This proves the lemma.\qedhere
\end{proof}

\medskip
\begin{lemma}[Circle actions and the slant product]\label{lem:sweeping-slant}
Let a circle act continuously on a space $Z$, with action map
$m:S^1\times Z\to Z$.  For a loop $a:S^1\to Z$, denote by
$\operatorname{sw}(a)$ the homology class
$m_*([S^1]\times[a])$.  If $\omega\in H^2(Z;\mathbb Z)$, define
$\omega/[S^1]=m^*\omega/[S^1]\in H^1(Z;\mathbb Z)$.  Then
$\langle\omega,\operatorname{sw}(a)\rangle
=\langle\omega/[S^1],[a]\rangle$.
\end{lemma}

\begin{proof}
Naturality of the Kronecker pairing and the defining adjunction for slant
product give
\[
 \begin{aligned}
 \bigl\langle\omega,\,
   m_*([S^1]\times[a])\bigr\rangle
 &=\bigl\langle m^*\omega,\,[S^1]\times[a]\bigr\rangle\\
 &=\bigl\langle m^*\omega/[S^1],\,[a]\bigr\rangle .
 \end{aligned}\qedhere
\]
\end{proof}

To compute the restriction to the cusp boundary, put
$\mathbf i=(i_1,i_2,i_3,i_4)
=\bigl(\gamma u^\vee,\gamma\Delta,
u^\vee\Delta-2\gamma w^\vee,Q\bigr)$.
This is a basis of $I_0$, since replacing $u^\vee w^\vee$ by
$Q=u^\vee w^\vee+3\gamma\Delta$ is a unimodular change of basis.
Choose lifts
$\widetilde{\mathbf i}=(\widetilde i_1,\ldots,\widetilde i_4)$ under
\eqref{eq:cusp-H2-splitting}.  Their restrictions to $M$ project to
$i_1,\ldots,i_4$ in the Wang sequence \eqref{eq:cusp-Wang}.  The two
solid-torus bounds in Lemma~\ref{lem:cusp-retraction} show that these
restrictions have zero $[w^\vee]$- and $[\Delta]$-coordinates.  After
adding $\kappa$ when necessary, they also have zero
$\vartheta$-component.  Write
$\bar i_k=\widetilde i_k|_M$, set
$\widetilde Q=\widetilde i_4$.  Thus
$(\bar i_1,\ldots,\bar i_4,[w^\vee],[\Delta])$ is a basis of
$H^2(M;\mathbb Z)/\operatorname{Tor}$.

The class $b$ may be replaced by $\pm b+\epsilon\alpha$, and
$Q_{\mathrm{fin}}$ by $Q_{\mathrm{fin}}+n b+\epsilon'\alpha$, where
$n\in\mathbb Z$ and $\epsilon,\epsilon'\in\{0,1\}$; these changes do not
alter their images in $H^2(F;\mathbb Z)$.  We now verify that they can be
chosen with boundary restrictions
\begin{equation}\label{eq:B-boundary-restrictions}
 b|_M=[w^\vee],\qquad
 \alpha|_M=\vartheta,\qquad
 Q_{\mathrm{fin}}|_M=4\widetilde Q|_M-[\Delta].
\end{equation}
Since $S_1,S_2$, and $F$ are aspherical and the attaching maps are
coverings, the standard tree-of-spaces universal cover of the homotopy
pushout is contractible.  Thus $Y_{\mathrm{fin}}$ is a $K(\pi,1)$.  We
use the standard description of classes in
$H^2(Y_{\mathrm{fin}};\mathbb Z)$ by central $\mathbb Z$-extensions.  To
find the free $C_0$-coordinates of $b|_M$,
we evaluate the corresponding central commutator on the $A_0$-invariant
fibre loops $w$ and $d$.  Under this description, the connecting class
$b=\delta[w^\vee]$ is obtained by gluing the trivial central extensions of
the two affine groups.  Let $\mathbf z$
denote the central generator, and let $\widetilde\lambda_j$ and
$\widetilde r_j$ be the standard lifts in the $j$-th extension.  Choose
the connecting-map convention
$\widetilde\lambda_1
=\mathbf z^{\,w^\vee(\lambda)}\widetilde\lambda_2$.
The lifts satisfy
$\widetilde r_j\widetilde\lambda_j\widetilde r_j^{-1}
=\widetilde{A_j\lambda}_j$.  Set
$\widetilde a_0=\widetilde r_2^{-1}\widetilde r_1^{-1}$; this lifts
$a_0=r_2^{-1}r_1^{-1}$.  Direct substitution gives
\[
 \widetilde a_0\widetilde\lambda_1\widetilde a_0^{-1}
 =\mathbf z^{\,w^\vee(A_1^{-1}\lambda)-w^\vee(A_0\lambda)}
   \widetilde{A_0\lambda}_1.
\]
Indeed, after conjugation by $\widetilde r_1^{-1}$, changing from the first
extension to the second contributes $w^\vee(A_1^{-1}\lambda)$;
conjugation by $\widetilde r_2^{-1}$ changes the translation to
$A_0\lambda$, and changing back contributes
$-w^\vee(A_0\lambda)$.  For an $A_0$-invariant $\lambda$, the central
commutator is therefore
$w^\vee(A_1^{-1}\lambda)-w^\vee(\lambda)$; reversing the convention changes
only its overall sign.  Since $A_1^{-1}w=-u+2d$ and $A_1^{-1}d=d$, its
image in the free part of $C_0$ is $\pm[w^\vee]$, with no
$[\Delta]$-component.  Changing the sign of $b$ and adding $\alpha$
gives the first formula in \eqref{eq:B-boundary-restrictions}.  The torsion
generator $d$ maps
isomorphically from $H_1(M)$ to the torsion summand of
$H_1(Y_{\mathrm{fin}})$, so
naturality of the universal-coefficient $\operatorname{Ext}$-term gives
the second formula.

Because $A_1d=A_2d=d$, the affine twists commute with translation by the
$d$-circle, which therefore acts globally
on $Y_{\mathrm{fin}}$.  We have oriented
$\operatorname{sw}_d(a_0)$ by $a_0\wedge d$.  In applying
Lemma~\ref{lem:sweeping-slant}, orient the acting circle by $-d$, so that
$(-d)\wedge a_0=a_0\wedge d$, and write slant product as
${-}/[-d]$.  Contraction on a smooth fibre then gives
$(4Q)/[-d]=12\gamma$.  Choose the generator
$\Theta\in H^1(Y_{\mathrm{fin}};\mathbb Z)$ normalized by
$\Theta(\mathfrak s)=1$.  Then $\Theta|_F=12\gamma$; injectivity of
$H^1(Y_{\mathrm{fin}})\to H^1(F)$ therefore gives
$Q_{\mathrm{fin}}/[-d]=\Theta$.  With the chosen orientation, the suspension torus
is the image of the product of the acting circle with the loop $a_0$, so
Lemma~\ref{lem:sweeping-slant} gives
$\langle Q_{\mathrm{fin}}|_M,\operatorname{sw}_{d}(a_0)\rangle
=\langle Q_{\mathrm{fin}}/[-d],a_0\rangle=\Theta(a_0)$.
Now $a_0=r_2^{-1}r_1^{-1}=-\mathfrak s$ in
$H_1(Y_{\mathrm{fin}})$, so $\Theta(a_0)=-1$.  The corresponding
suspension torus bounds in $A$, and hence pairs trivially with
$4\widetilde Q|_M$.  Consequently,
$\langle Q_{\mathrm{fin}}|_M-4\widetilde Q|_M,
\operatorname{sw}_{d}(a_0)\rangle=-1$.
This is the coefficient of $[\Delta]$ in
$Q_{\mathrm{fin}}|_M-4\widetilde Q|_M$, with the stated sign.

For the Mayer--Vietoris difference map we use the convention
$\operatorname{res}_A-\operatorname{res}_{Y_{\mathrm{fin}}}$.
In the domain basis
$(\widetilde{\mathbf i},\kappa,b,Q_{\mathrm{fin}})$ and the target basis
$(\bar i_1,\bar i_2,\bar i_3,\bar i_4,[w^\vee],[\Delta])$ of the free
parts, its matrix is
\begin{equation}\label{eq:boundary-free-matrix}
 \begin{pmatrix}
 1&0&0&0&0& 0& 0\\
 0&1&0&0&0& 0& 0\\
 0&0&1&0&0& 0& 0\\
 0&0&0&1&0& 0&-4\\
 0&0&0&0&0&-1& 0\\
 0&0&0&0&0& 0& 1
 \end{pmatrix}.
\end{equation}
Its six nonzero Smith invariants are all one.  The remaining torsion row is
the map
\begin{equation}\label{eq:boundary-torsion-row}
 \mathbb Z\langle\kappa\rangle\oplus
 (\mathbb Z/2)\langle\alpha\rangle
 \longrightarrow(\mathbb Z/2)\langle\vartheta\rangle,
 \qquad (n,\epsilon)\longmapsto n-\epsilon\pmod2.
\end{equation}
Equations \eqref{eq:boundary-free-matrix} and
\eqref{eq:boundary-torsion-row} are the free and torsion parts of the same
integral map.

\subsection{The Mayer--Vietoris map in cohomological degree two}

The cusp neighbourhood and the neighbourhood of the multiplicity-four fibre
contribute classes whose restrictions to the common boundary have order two.
Equations~\eqref{eq:kappa-boundary}
and \eqref{eq:B-boundary-restrictions} show that their restrictions are equal.

\begin{proof}[Proof of Proposition~\ref{prop:degree-two-mv}]
The classes $\kappa$ and $\alpha$ have the same nonzero restriction,
$\kappa|_M=\vartheta=\alpha|_M$,
by \eqref{eq:kappa-boundary} and
\eqref{eq:B-boundary-restrictions}.  To determine the full kernel, use
\eqref{eq:cusp-H2-splitting}, \eqref{eq:cusp-Wang}, and
\eqref{eq:B-boundary-restrictions}.  The $I_0$-part first cancels the
$4Q$-part of a multiple of $Q_{\mathrm{fin}}$.  Its remaining coordinate in
$C_0$ is the same multiple of the primitive class $-[\Delta]$, so
that multiple is zero.  The primitive $[w^\vee]$-coordinate then excludes the
coefficient of the connecting generator $b$.  The remaining pairs are
$(n\kappa,\epsilon\alpha)$, where $n\in\mathbb Z$,
$\epsilon\in\{0,1\}$, and $n\equiv\epsilon\pmod2$.
They form the infinite cyclic group generated by
\eqref{eq:x-generator}.  Since $2\alpha=0$,
\eqref{eq:cusp-component-class} becomes
$[D_0]=(2\kappa,0)=2x$ and $[D_1]=-2x$.
\end{proof}

\subsection{Van Kampen and Mayer--Vietoris}

\begin{proof}[Proof of Proposition~\ref{prop:integral-topology}]
Let $r_1,r_2$ be the affine generators.  Lemma~\ref{lem:cusp-retraction}
and van Kampen give
\[
 \begin{split}
 \pi_1(Y)=\langle\Lambda^+,r_1,r_2\mid{}&
 r_j\lambda r_j^{-1}=A_j\lambda,
 \quad r_1^3=v_1,\quad r_2^4=v_2,\\
 &w=d=1,\quad r_1r_2=1\rangle.
 \end{split}
\]
The normal closure of $w,d$ also kills $u$, for
$A_2w=-u+2d$.  If $c$ is the image of $\widehat\gamma$, and
$x_1=r_1,x_2=r_2$, the remaining relations are
$x_1^3=c$, $x_2^4=c^{-1}$, and $x_1x_2=1$.
Thus $c=x_1^3=x_1^4$, and every generator is trivial.  Hence
$\pi_1(Y)=1$.

The map on first homology retains the integral torsion.  From
\eqref{eq:cusp-monodromy-homology},
$H_1(M)=\mathbb Z\langle\widehat\gamma,u,a_0\rangle
\oplus(\mathbb Z/2)\langle d\rangle$.
The map to the two pieces is
\begin{equation*}
 \begin{array}{c|cccc}
  &\widehat\gamma&u&a_0&d\\ \hline
 H_1(A)&\widehat\gamma&u&0&0\\
 H_1(Y_{\mathrm{fin}})&12\mathfrak s&0&-\mathfrak s&d.
 \end{array}
\end{equation*}
Its determinant on the free summands is $-1$, and it is the identity on
the summand of order two.  It is therefore an isomorphism.

Cohomological Mayer--Vietoris consequently identifies $H^2(Y;\mathbb Z)$
with the kernel in
Proposition~\ref{prop:degree-two-mv}.  Hence
$H^2(Y;\mathbb Z)=\mathbb Zx$ and
$[D_0]=2x$, $[D_1]=-2x$.

The same boundary formulas show that
$H^2(A)\oplus H^2(Y_{\mathrm{fin}})\to H^2(M)$ is surjective.  Its
composition with $H^2(M)\to I_0$ is surjective, and its image also contains
$\vartheta$ and $[w^\vee]$.  The difference of
$4\widetilde Q$ and $Q_{\mathrm{fin}}$ gives the primitive class
$[\Delta]$.  Mayer--Vietoris therefore injects
\begin{equation}\label{eq:H3-Y-injection}
 H^3(Y;\mathbb Z)\hookrightarrow
 H^3(A;\mathbb Z)\oplus H^3(Y_{\mathrm{fin}};\mathbb Z).
\end{equation}
The target is torsion free by \eqref{eq:cusp-homology} and
\eqref{eq:H3-B-torsion-free}.

For the Euler characteristic, use $e_c$, which is
additive for this locally closed stratification of the base and
multiplicative for the locally trivial part of the fibration.  Since $Y$
is compact, $e(Y)=e_c(Y)$.  Each reduced multiple fibre is a free finite
quotient of a real four-torus, so its Euler characteristic is zero; the
scheme-theoretic multiplicity does not change its underlying topological
space.  Over the remaining punctured-base stratum the fibre is a real
four-torus, and hence
$e_c(f^{-1}(B^\circ))=e_c(B^\circ)e(T^4)=0$.
The cusp fibre is the only contribution, and
\eqref{eq:cusp-homology} gives $e(W)=e(A)=4$.  Therefore $e(Y)=4$.
Simple connectivity, $H^2(Y)=\mathbb Z$, Poincar\'e duality, and this
Euler number give $b_3(Y)=0$.  The injection
\eqref{eq:H3-Y-injection} then gives $H^3(Y;\mathbb Z)=0$, including
torsion.  The injection
$\operatorname{Ext}(H_2(Y;\mathbb Z),\mathbb Z)\hookrightarrow
H^3(Y;\mathbb Z)=0$ shows that $H_2(Y;\mathbb Z)$ is torsion free,
and integral Poincar\'e duality supplies all remaining groups.
\end{proof}


\section{Characteristic classes}
\label{sec:characteristic-classes}

Let $D$ be either component of the cusp fibre, and choose the sign of
$x\in H^2(Y;\mathbb Z)$ so that $[D]=2x$.  The normalization of $D$
and its conductor determine the cubic form and the characteristic classes.

\begin{proof}[Proof of Proposition~\ref{prop:characteristic-classes}]
Let $\nu:S\to D$, $C_1,\ldots,C_6$, and
$C_{\mathrm{dbl}}$ be as in Lemma~\ref{lem:cusp-double-curve}.  That lemma
gives $D^3=0$, $D\cdot C_{\mathrm{dbl}}=-2$, and
$\deg N_{C_{\mathrm{dbl}}/Y}=-2$.
Thus $8x^3=0$; the group $H^6(Y;\mathbb Z)\cong\mathbb Z$ is
torsion free, so $x^3=0$.  Moreover,
$[D]=2x$ gives $x(C_{\mathrm{dbl}})=-1$, and hence
$C_{\mathrm{dbl}}$ is a primitive generator of
$H_2(Y;\mathbb Z)$.  Adjunction on this smooth rational curve gives
$\langle c_1(Y),C_{\mathrm{dbl}}\rangle
=\deg T\mathbb P^1+\deg N_{C_{\mathrm{dbl}}/Y}=2-2=0$.
It follows that
\begin{equation}\label{eq:c1-zero}
 c_1(Y)=0,\qquad w_2(Y)=0.
\end{equation}
Here $c_1(Y)=0$ is an equality of integral topological Chern classes; it
does not assert a holomorphic trivialization of $K_Y$.

Let $\phi:C_1\to C_4$ be the isomorphism that identifies the two branches.
The normalization sequence for $D$ is
\begin{equation}\label{eq:normalization-sequence}
 0\longrightarrow\mathcal O_D
 \longrightarrow\nu_*\mathcal O_S
 \xrightarrow{\operatorname{res}_{C_1}
  -\phi^*\operatorname{res}_{C_4}}
 \mathcal O_{C_{\mathrm{dbl}}}\longrightarrow0.
\end{equation}
At a triple point the local algebra of $D$ and its normalization is
\[
 \frac{\mathbb C\{z_1,z_2,z_3\}}{(z_1z_2)}
 \longrightarrow
 \mathbb C\{z_1,z_3\}\oplus\mathbb C\{z_2,z_3\}.
\]
The cokernel is $\mathbb C\{z_3\}$, via the difference of the two
restrictions at $z_1=z_2=0$.  Hence the difference map in
\eqref{eq:normalization-sequence} remains surjective at the triple point and
has no additional zero-dimensional cokernel.  Since $S$ is
rational and $C_{\mathrm{dbl}}\cong\mathbb P^1$, equation
\eqref{eq:normalization-sequence} gives $\chi(\mathcal O_D)=1-1=0$.
On the smooth threefold $Y$, apply the Hirzebruch--Riemann--Roch theorem
to the $K$-theory identity supplied by
$0\to\mathcal O_Y(-D)\to\mathcal O_Y\to\mathcal O_D\to0$.
Thus $[\mathcal O_D]=1-[\mathcal O_Y(-D)]$ and
$\operatorname{ch}(\mathcal O_D)=1-e^{-D}$.  Using
\eqref{eq:c1-zero}, the degree-six part of Hirzebruch--Riemann--Roch is
$\chi(\mathcal O_D)
=\int_Y(D^3/6+D\,c_2(Y)/12)$.
Since $D^3=\chi(\mathcal O_D)=0$, the Riemann--Roch formula gives
$D\,c_2(Y)=0$.  Choose $y\in H^4(Y;\mathbb Z)$ so that
$\langle x\smile y,[Y]\rangle=1$, and write
$c_2(Y)=ay$.  Since $D=2x$, we obtain $2a=0$, and hence
$c_2(Y)=0$ and $p_1(Y)=c_1(Y)^2-2c_2(Y)=0$.
\end{proof}
